\documentclass[a4paper,12pt]{article}
\usepackage[utf8]{inputenc}
\usepackage[portuguese]{babel}
\usepackage[T1]{fontenc}
\usepackage{geometry}
\geometry{left=2.5cm, right=2.5cm, top=2.5cm, bottom=2.5cm}
\usepackage{graphicx}
\usepackage{xcolor}
\usepackage{tikz}
\usetikzlibrary{shapes,arrows,positioning,calc}
\usepackage{tcolorbox}
\tcbuselibrary{skins,breakable}
\usepackage{hyperref}
\hypersetup{colorlinks=true, linkcolor=blue, urlcolor=blue}
\usepackage{fancyhdr}
\usepackage{lastpage}
\usepackage{fontawesome5} % Para ícones (requer pacote fontawesome5 ou similar)

% Configuração de cabeçalho e rodapé
\pagestyle{fancy}
\fancyhf{}
\fancyhead[L]{\small\textit{Guia de Negócio: Digitalização de Restaurantes}}
\fancyhead[R]{\small\textit{Kilamba}}
\fancyfoot[C]{\thepage\ de \pageref{LastPage}}
\renewcommand{\headrulewidth}{0.4pt}
\renewcommand{\footrulewidth}{0.4pt}

% Cores personalizadas
\definecolor{corprincipal}{RGB}{0,102,204}
\definecolor{cordestaque}{RGB}{204,102,0}
\definecolor{corfundo}{RGB}{245,245,245}

% Comando para caixas de dica
\newtcolorbox{dica}{
    colback=corfundo,
    colframe=corprincipal,
    title=Dica,
    fonttitle=\bfseries,
    boxrule=0.5mm,
    arc=3mm,
    breakable
}

% Comando para caixas de alerta
\newtcolorbox{alerta}{
    colback=red!5,
    colframe=red!75!black,
    title=Atenção,
    fonttitle=\bfseries,
    boxrule=0.5mm,
    arc=3mm,
    breakable
}

% Comando para caixas de sucesso
\newtcolorbox{sucesso}{
    colback=green!5,
    colframe=green!75!black,
    title=Sucesso,
    fonttitle=\bfseries,
    boxrule=0.5mm,
    arc=3mm,
    breakable
}

\begin{document}

\begin{center}
    {\Huge\bfseries\color{corprincipal} Guia Completo para o Sucesso}\\[0.3cm]
    {\Large\bfseries Negócio de Digitalização de Restaurantes em Kilamba}\\[1cm]
    {\large \textit{Um passo a passo estratégico para capturar o mercado local}}\\[0.5cm]
    \rule{\textwidth}{0.4pt}\\[0.5cm]
    {\large \textbf{Data:} \today}\\[0.2cm]
    {\large \textbf{Versão:} 1.0}\\[0.2cm]
    \vspace{0.5cm}
    \includegraphics[width=0.2\textwidth]{example-image} % Substitua por um logo se desejar
\end{center}

\tableofcontents
\newpage

\section{Introdução}
O mercado de restaurantes em Kilamba apresenta uma oportunidade única para serviços de marketing digital e gestão de presença online. Com o crescimento do uso do Google Maps e a procura por comodidade, os estabelecimentos que não possuem um perfil digital bem estruturado perdem clientes diariamente.

Este documento fornece um roteiro detalhado para criar um negócio de sucesso que oferece:
\begin{itemize}
    \item Otimização de perfis no Google Maps (Google Business Profile);
    \item Criação de cardápios digitais atualizáveis via QR Code;
    \item Gestão de reputação online e análise de tendências com Google Trends;
    \item Um pacote completo de serviços que aumenta a visibilidade e as vendas dos restaurantes.
\end{itemize}

Siga cada fase para construir uma base sólida, desde a pesquisa inicial até a venda e implementação.

\section{Fase 1: Pesquisa e Análise de Mercado}
Antes de abordar qualquer restaurante, é essencial entender o cenário local. Utilize as seguintes etapas:

\subsection{Mapear os Spots Mais Populares}
\begin{enumerate}
    \item Acesse o Google Maps e pesquise por “restaurantes em Kilamba”, “melhores spots Kilamba”, “onde comer em Kilamba”.
    \item Filtre pelos estabelecimentos com maior número de avaliações (acima de 100) e pontuação média superior a 4,0.
    \item Crie uma lista com os 10–15 mais populares, anotando:
    \begin{itemize}
        \item Nome e endereço;
        \item Se possuem perfil no Google Maps completo (fotos, horários, website);
        \item Tipo de cozinha (hambúrguer, pizza, comida fancy, etc.);
        \item Presença em redes sociais e site próprio.
    \end{itemize}
    \item \textbf{Objetivo:} Identificar os concorrentes diretos do seu futuro cliente e entender o que já funciona bem.
\end{enumerate}

\subsection{Analisar Outros Spots na Mesma Área}
\begin{enumerate}
    \item Navegue manualmente pelo mapa de Kilamba no Google Maps, zoomando em ruas e bairros.
    \item Procure restaurantes que tenham poucas avaliações (menos de 50) ou informações incompletas.
    \item Esses serão seus \textbf{clientes em potencial} — eles estão “invisíveis” digitalmente e carecem de serviços de marketing.
    \item Registre contatos (telefone, Instagram) para abordagem posterior.
\end{enumerate}

\subsection{Usar o Google Trends para Validar a Demanda}
O Google Trends ajuda a provar que as pessoas já estão procurando ativamente por comidas específicas em Kilamba.
\begin{enumerate}
    \item Acesse \url{https://trends.google.com}.
    \item Compare termos como:
    \begin{itemize}
        \item “pizza Kilamba”
        \item “hambúrguer Kilamba”
        \item “comida fancy Kilamba”
        \item “restaurante familiar Kilamba”
    \end{itemize}
    \item Analise os gráficos de tendência (picos em fins de semana, feriados).
    \item \textbf{Exemplo de argumento:} “Olha, as pesquisas por ‘pizza Kilamba’ aumentam 80\% aos sábados, mas seu restaurante não aparece nos resultados. Com meu serviço, você estará visível exatamente quando o cliente está procurando.”
\end{enumerate}

\begin{dica}
    Capture prints dos gráficos do Trends e monte um pequeno relatório para apresentar aos clientes. Isso demonstra profissionalismo e dados concretos.
\end{dica}

\section{Fase 2: Ferramentas e Preparação do Pacote de Serviços}
Agora você definirá os serviços que oferecerá, as ferramentas utilizadas e como estruturar o pacote.

\subsection{Otimização do Google Maps (Google Business Profile)}
Cada restaurante precisa de um perfil completo e ativo no Google. O serviço inclui:
\begin{itemize}
    \item \textbf{Criação/Reivindicação:} Se o estabelecimento não tem perfil, criar; se tem, reivindicar a propriedade.
    \item \textbf{Informações completas:} Horários, telefone, website, categoria correta.
    \item \textbf{Fotos profissionais:} Fotografar pratos, ambiente, fachada (ou contratar um fotógrafo).
    \item \textbf{Atributos:} Marcar opções como “tem estacionamento”, “aceita reservas”, “tem música ao vivo”, etc.
    \item \textbf{Gestão de avaliações:} Responder a comentários positivos e negativos de forma cordial.
\end{itemize}

\subsection{Criação do Cardápio Digital}
A principal inovação que você venderá é o cardápio digital atualizável. Duas abordagens principais:

\subsubsection{Cardápio via QR Code (Sistemas especializados)}
\begin{itemize}
    \item \textbf{Ferramentas recomendadas:} OlaClick, QR Menu Maker, MenuQ.
    \item \textbf{Vantagens:}
    \begin{itemize}
        \item O cliente escaneia o código na mesa e vê o cardápio no celular.
        \item Atualizações em tempo real (preços, disponibilidade).
        \item Permite incluir fotos de alta qualidade, vídeos e até pedidos online (em versões mais avançadas).
    \end{itemize}
    \item \textbf{Custo:} A maioria oferece planos mensais acessíveis (5–20 USD/mês). Você pode repassar esse custo ao cliente como parte do pacote.
\end{itemize}

\subsubsection{Cardápio no Site ou WordPress}
\begin{itemize}
    \item \textbf{Ferramentas:} WordPress com plugins como MenuMax, ou construtores como Wix, Canva (para PDF interativo).
    \item \textbf{Vantagem:} Cria uma presença institucional mais robusta; pode ser integrado ao Google Maps.
    \item \textbf{Desvantagem:} Atualizações podem exigir mais tempo se não forem feitas via sistema simplificado.
\end{itemize}

\begin{sucesso}
    O case da OlaClick mostra aumento de 25\% nas vendas de restaurantes que adotaram o cardápio digital com QR Code. Utilize esse dado para convencer os clientes.
\end{sucesso}

\subsection{Design dos QR Codes}
Não basta gerar o código; é necessário apresentá-lo de forma atraente.
\begin{itemize}
    \item Use o \textbf{Canva} ou \textbf{Adobe Express} para criar templates personalizados com o logo do restaurante.
    \item Ofereça opções de suportes: adesivos para mesas, porta-códigos de acrílico, ou até cardápios físicos reduzidos com o QR Code destacado.
\end{itemize}

\section{Fase 3: O Discurso de Vendas (O Pitch)}
Um bom pitch é a diferença entre conquistar ou perder um cliente. Estruture sua conversa em três momentos:

\subsection{Diagnóstico Personalizado}
\begin{quote}
    “Bom dia, \[nome do proprietário\]. Passei pelo seu restaurante e notei que, no Google Maps, ele não aparece quando alguém procura por \[tipo de comida\] aqui em Kilamba. Além disso, o cardápio atual parece ser apenas impresso, o que dificulta mudanças rápidas. Estou aqui para apresentar uma solução que resolve esses dois problemas e ainda aumenta as vendas.”
\end{quote}

\subsection{A Solução (O Pacote Completo)}
Apresente os três pilares do seu serviço:
\begin{enumerate}
    \item \textbf{Visibilidade no Google:} Criação/otimização do perfil, fotos profissionais, gestão de avaliações.
    \item \textbf{Cardápio Digital:} Sistema com QR Code, atualizável em tempo real, sem custos de impressão.
    \item \textbf{Acompanhamento:} Relatórios mensais de desempenho (quantas visualizações no Google, cliques no cardápio).
\end{enumerate}

\subsection{Demonstração de Valor (Por que isso dá dinheiro?)}
Utilize dados:
\begin{itemize}
    \item “Sabia que 77\% dos consumidores consultam o menu online antes de decidir onde comer? Sem presença digital, você perde esses clientes.”
    \item “Com o cardápio digital, o ticket médio costuma aumentar porque fotos apetitosas incentivam pedidos adicionais.”
    \item “Restaurantes que respondem a avaliações no Google têm 30\% mais chances de serem escolhidos.”
\end{itemize}

\begin{dica}
    Tenha sempre um \textbf{demo} pronto. Use um dispositivo móvel para mostrar como funciona o QR Code: escaneie, abra o cardápio, simule uma atualização de preço em tempo real. Isso gera impacto imediato.
\end{dica}

\section{Fase 4: Implementação e Manutenção}
Depois de fechar o contrato, siga este fluxo de trabalho:

\subsection{Criação da Demonstração (Demo) Personalizada}
\begin{enumerate}
    \item Selecione um restaurante piloto ou crie um exemplo fictício.
    \item Configure um cardápio digital básico com alguns pratos.
    \item Mostre ao cliente o painel de administração: como adicionar um novo prato, alterar preços, desativar itens.
    \item Entregue uma amostra física do QR Code para que ele possa testar.
\end{enumerate}

\subsection{Impressão e Instalação dos QR Codes}
\begin{itemize}
    \item Utilize papel adesivo resistente ou suportes de acrílico para colocar nas mesas.
    \item Coloque também um QR Code na entrada (para clientes que passam na rua) e outro no balcão de pedidos (se houver).
    \item Recomende que o restaurante mantenha um cardápio físico reduzido para os clientes que preferem o tradicional, mas sempre com destaque para o código.
\end{itemize}

\subsection{Estrutura de Preços}
Ofereça modelos de assinatura mensal para garantir receita recorrente. Sugestão:

\vspace{0.5cm}
\begin{center}
\begin{tabular}{|p{3cm}|p{7cm}|p{4cm}|}
\hline
\multicolumn{1}{|c|}{\textbf{Plano}} & \multicolumn{1}{c|}{\textbf{O que inclui}} & \multicolumn{1}{c|}{\textbf{Preço sugerido}} \\
\hline
\textbf{Básico} & 
- Criação/otimização do perfil Google Business Profile \\
- Cardápio digital com QR Code (até 30 itens) \\
- 1 atualização de cardápio por mês \\
- Suporte por e-mail &
AOA 15.000–25.000/mês \\
\hline
\textbf{Premium} & 
- Tudo do Básico \\
- Atualizações ilimitadas \\
- Gestão de avaliações no Google (respostas a todos os comentários) \\
- Postagens mensais no perfil do Google (promoções) \\
- Relatório mensal de desempenho &
AOA 35.000–50.000/mês \\
\hline
\textbf{Setup Único} & 
- Fotografia profissional dos pratos e ambiente \\
- Design personalizado do QR Code \\
- Configuração inicial completa (Google + cardápio) &
AOA 30.000–60.000 (uma vez) \\
\hline
\end{tabular}
\end{center}
\vspace{0.5cm}

\subsection{Fluxo de Trabalho (Workflow)}
Abaixo um diagrama visual do processo completo:

\begin{center}
\begin{tikzpicture}[node distance=1.2cm, auto, font=\small,
    box/.style={rectangle, draw=corprincipal, thick, fill=corfundo, text width=3.5cm, align=center, minimum height=0.8cm, rounded corners=2mm},
    arrow/.style={->, >=stealth, thick}]
    \node[box] (pesquisa) {Pesquisa e Mapeamento};
    \node[box, below=of pesquisa] (contato) {Abordagem e Pitch};
    \node[box, below=of contato] (contrato) {Fechamento do Contrato};
    \node[box, below=of contrato] (setup) {Setup: Google + Cardápio Digital};
    \node[box, below=of setup] (entrega) {Entrega: QR Codes + Treinamento};
    \node[box, below=of entrega] (manutencao) {Manutenção Mensal};
    
    \draw[arrow] (pesquisa) -- (contato);
    \draw[arrow] (contato) -- (contrato);
    \draw[arrow] (contrato) -- (setup);
    \draw[arrow] (setup) -- (entrega);
    \draw[arrow] (entrega) -- (manutencao);
    \draw[arrow] (manutencao.west) -- ++(-1,0) |- (pesquisa.west);
\end{tikzpicture}
\end{center}

\section{Conclusão e Próximos Passos}
Com este guia, você tem em mãos uma metodologia testada para transformar a presença digital de restaurantes em Kilamba. Os próximos passos práticos são:
\begin{enumerate}
    \item Realizar a pesquisa de mercado nos próximos 3 dias.
    \item Criar um portfólio de demonstração (exemplo fictício) com um cardápio digital pronto.
    \item Definir os preços finais com base na sua região e público-alvo.
    \item Agendar reuniões com os 10 primeiros potenciais clientes mapeados.
    \item Começar com um piloto (um cliente que aceite um desconto em troca de um depoimento) e usar esse case para alavancar novos contratos.
\end{enumerate}

Lembre-se: a chave do sucesso é \textbf{mostrar resultados}. Cada novo cliente deve gerar um aumento mensurável em visualizações, cliques e, consequentemente, vendas. Utilize esses números para refinar seu discurso e crescer de forma escalável.

\begin{center}
    \rule{0.7\textwidth}{0.4pt}\\
    \textit{“O melhor marketing não é aquele que gasta mais, mas aquele que está presente quando o cliente decide.”}\\
    \textbf{Boa sorte e bons negócios!}
\end{center}

\end{document}